\chapter{The Bifrost WORM Chain Protocol}
\label{ch:bifrost}

\section{Purpose}
\label{sec:bifrost-purpose}

The Bifrost WORM (Write-Once-Read-Many) Chain is the trust root for
MATHLIB5. It is an append-only hash chain in which each block anchors a
verification receipt (or a batch of receipts) produced by a pipeline
boundary. Because the chain is write-once and the head is anchored to a
Bitcoin \texttt{OP\_RETURN} transaction (\cref{ch:anchors}), the provenance
of every artifact becomes externally, economically, auditable.

\section{Block Structure}
\label{sec:bifrost-block}

Each Bifrost block is the tuple
\[
\mathcal{B}_n = \bigl(n,\; h_{n-1},\; \mathsf{merkle}(\mathcal{R}_n),\;
t_n,\; \sigma_n \bigr),
\]
where:
\begin{itemize}
  \item $n$ is the block height,
  \item $h_{n-1} = \text{SHA-256}(\mathcal{B}_{n-1})$ is the previous block
        hash (the ``worm'' link),
  \item $\mathcal{R}_n$ is the set of receipts anchored at this height,
  \item $\mathsf{merkle}(\mathcal{R}_n)$ is a Merkle root binding all
        receipts,
  \item $t_n$ is the timestamp,
  \item $\sigma_n$ is an Ed25519 signature by the Plasma Gate over the
        above fields.
\end{itemize}
The genesis block $\mathcal{B}_0$ carries the fixed trust-root identifier
\texttt{4b565498-9afc-4782-af4a-c6b11a5d0058}.

\section{Anchor Format and the Bitcoin Link}
\label{sec:bifrost-anchor}

To publish a Bifrost head externally, MATHLIB5 emits a Bitcoin
\texttt{OP\_RETURN} output encoding
\[
\texttt{MATH5:}\;\Vert\; \text{SHA-256}(\mathcal{B}_{\text{head}}).
\]
The \texttt{MATH5:} prefix (6 bytes) namespaces the message; the 32-byte
hash follows, for a total of 38 payload bytes (plus the 1-byte push opcode
and 2-byte \texttt{OP\_RETURN}, well within the 80-byte standard limit).
A verifier on any machine can:
\begin{enumerate}
  \item Fetch the Bitcoin transaction from the public ledger.
  \item Recompute $\text{SHA-256}(\mathcal{B}_{\text{head}})$ locally.
  \item Confirm \texttt{MATH5:} matches and the hashes agree.
\end{enumerate}
If they agree, the local artifact chain is provably consistent with the
publicly anchored one as of that transaction.

\section{Inclusion Proofs}
\label{sec:bifrost-inclusion}

A single receipt $r \in \mathcal{R}_n$ is proven included by a Merkle
path $m_r$ such that
\[
\mathsf{merkle}(\mathcal{R}_n) \;=\; \mathsf{reduce}(m_r,\, \text{SHA-256}(r)).
\]
Combined with the chain hash link, this yields a \emph{succinct} proof that
$r$ was anchored at height $n$ without revealing the other receipts.

\section{Trust Properties}
\label{sec:bifrost-trust}

\begin{itemize}
  \item \textbf{Immutability.} Tampering with any block $k \le n$ changes
        $h_k$, which cascades to change $h_n$ and therefore the Bitcoin
        anchor --- detectable by anyone.
  \item \textbf{External anchoring.} The Bitcoin network's proof-of-work
        makes rewriting history costly in proportion to its depth, lifting
        the trust root out of MATHLIB5's own control.
  \item \textbf{Seal-and-verify.} Every boundary receipt is independently
        signable, so the chain is a fine-grained audit trail, not a coarse
        commit log.
\end{itemize}
These properties realize the AGENTS.md claim that ``the repo \emph{is} the
training curve'': each verified solution is a new immutable memory bucket,
anchored in the same chain that secures the toolchain itself.
